\documentclass[12pt]{article}
\usepackage[utf8]{inputenc}
\usepackage[T1]{fontenc}
\usepackage[margin=1in]{geometry}
\usepackage{setspace}
\usepackage{xcolor}
\usepackage[colorlinks=true, urlcolor=blue!60!black, linkcolor=blue!60!black]{hyperref}
\setlength{\parindent}{2em}
\setlength{\parskip}{1.25em}
\renewcommand{\baselinestretch}{1.0}
\begin{document}
\begin{center}
{\LARGE \textit{Home Building \& Loan Ass'n v. Blaisdell}, 290 U.S. 398 (1934)\par}
\vspace{0.8em}
{\large Case Brief --- Four-Page Edition\par}
\end{center}
\vspace{1em}

\section{Memory Jogger}

During the Great Depression, Minnesota let courts temporarily extend mortgage-redemption periods if the mortgagor paid a court-set rental value; a 5--4 Court upheld the law, holding that the Contracts Clause is not absolute and yields to a State's reasonable, temporary, conditioned exercise of its police power in an emergency (\href{https://supreme.justia.com/cases/federal/us/290/398/}{Justia}).

\section{Detailed Case Facts}

As the Depression drove property values down and foreclosures up, Minnesota enacted the Mortgage Moratorium Law on April 18, 1933, declaring a public economic emergency. The law did not cancel debts or impose a blanket freeze. It let state courts, on a mortgagor's application, extend the redemption period after a foreclosure sale as was just and equitable, but no later than May 1, 1935, on condition that the mortgagor in possession pay all or part of the property's reasonable rental value toward taxes, insurance, interest, and the debt. Courts could shorten extensions if circumstances changed. The Minneapolis home of John H.\ Blaisdell and his wife was foreclosed on May 2, 1932, and bought by the mortgagee, Home Building \& Loan Association; their one-year redemption period would expire May 2, 1933.

\section{Procedural History}

The Blaisdells petitioned for an extension; the Association argued that the law violated the Contracts Clause and the Fourteenth Amendment. The trial court held the law unconstitutional, but the Minnesota Supreme Court reversed. On remand, the trial court set the rental value at \$40 a month and extended redemption to May 1, 1935, conditioned on monthly payments. The Minnesota Supreme Court affirmed, and the Association appealed. The U.S. Supreme Court affirmed on January 8, 1934.

\section{Judicial Votes}

Affirmed, \textbf{5--4}.
\begin{itemize}
\item \textbf{Majority (5):} Chief Justice Hughes, joined by Justices Brandeis, Stone, Roberts, and Cardozo.
\item \textbf{Dissent (4):} Justice Sutherland, joined by Justices Van Devanter, McReynolds, and Butler.
\end{itemize}

\section{Holding}

The Contracts Clause, which bars any State from passing a ``Law impairing the Obligation of Contracts,'' is not an absolute prohibition applied with mathematical exactness; it must be harmonized with the State's reserved police power to protect its people's vital interests in an emergency. Because Minnesota's relief was temporary, limited to a genuine emergency, directed to a legitimate public end, and conditioned on reasonable terms protecting the mortgagee's investment, it was valid. For the same reasons, it satisfied due process, and its classification did not deny equal protection.

\section{Analysis of Opinions}

\textbf{Majority (Hughes).} An emergency neither creates nor enlarges power, but it may furnish the occasion for exercising power the State already has. The Clause arose to curb post-Revolutionary debtor-relief laws, yet it is a broad provision requiring construction, and every contract is made subject to the State's continuing authority (\textit{Stone v. Mississippi}; \textit{Block v. Hirsh}). Rejecting formal distinctions between right and remedy, Hughes asked ``whether the legislation is addressed to a legitimate end and the measures taken are reasonable and appropriate to that end.'' Unlike the unconditional extensions in \textit{Bronson v. Kinzie} and \textit{Barnitz v. Beverly}, Minnesota compensated the mortgagee. Rejecting static originalism, he read the ``great clauses of the Constitution'' in light of the Nation's ``whole experience.'' Five factors saved the law: a genuine emergency, a public purpose, tailored judicial relief, protection of the mortgagee's core interest (the debt, interest, sale, and rent all remained), and temporary operation, with the emergency's continuance ``always open to judicial inquiry.''

\textbf{Sutherland, with Van Devanter, McReynolds, and Butler (dissenting).} A constitutional provision's meaning does not change with economic conditions. The Clause was drafted precisely to bar debtor-relief moratoria in depressions, so upholding one because of a depression turns it upside down. Under \textit{Bronson} and \textit{Barnitz}, extending redemption deprives the purchaser of bargained-for title and possession, and court-set rent is no substitute. An open-ended ``emergency'' balancing test invites erosion of constitutional limits.

\textbf{Later developments.} The Court struck permanent or uncompensated relief in \textit{W. B. Worthen Co. v. Thomas} and \textit{Kavanaugh}, then dropped the emergency requirement in \textit{Veix} (1940). \textit{United States Trust} (1977) and \textit{Allied Structural Steel} (1978) briefly revived the Clause; \textit{Energy Reserves} (1983) set a deferential three-part test; and \textit{Sveen v. Melin} (2018) streamlined it, with Justice Gorsuch dissenting.

\section{Examples: Future Applications}

\begin{enumerate}
\item \textbf{Same-Side: Compensated Pandemic Foreclosure Pause.} A State lets courts postpone foreclosure sales up to twelve months for homeowners who lost jobs in a declared emergency, if they pay fair rental value; the law expires with the emergency. Every \textit{Blaisdell} factor favors the State.
\item \textbf{Same-Side: Notice Rule in a Regulated Industry.} A State requires solar lessors to mail a one-page notice before automatic renewals. Heavy prior regulation and minimal paperwork mean no substantial impairment under \textit{Sveen}.
\item \textbf{Opposite-Side: Permanent Rent Forgiveness.} After a hurricane, a State permanently cancels six months' commercial rent and guaranties without hardship findings or compensation, like \textit{Kavanaugh}.
\item \textbf{Opposite-Side: Unconditional Five-Year Extension.} Without an acute crisis, a State extends every redemption period five years, with no hardship showing, rent, court supervision, or expiration. It lacks every guardrail that saved Minnesota's law.
\item \textbf{Fence-Sitter: Repealing Retirees' COLA.} A State repeals a 3\% cost-of-living adjustment for retired teachers. \textit{U.S. Trust} heightened scrutiny applies to public contracts, but a court must first decide whether the statute created a contract at all.
\end{enumerate}

\section{Critique}

\textit{Blaisdell} pits living constitutionalism against originalism on a clause expressly aimed at debtor-relief stay laws. Law-and-economics critics warn that ex post relief raises ex ante credit costs, and the case shows how ``emergency'' exceptions harden into routine deference. Because Article I, \S\ 10 binds only States, federal retroactivity is reviewed instead under Fifth Amendment due process or the Takings Clause.

\section{Key Quotations}

\begin{quote}
The Contracts Clause's prohibition ``is not an absolute one, and is not to be read with literal exactness, like a mathematical formula.''
\end{quote}
\noindent --- Chief Justice Hughes, 290 U.S. at 428.

\begin{quote}
``Emergency does not create power. Emergency does not increase granted power or remove or diminish the restrictions imposed upon power granted or reserved. The Constitution was adopted in a period of grave emergency. Its grants of power to the federal government and its limitations of the power of the States were determined in the light of emergency, and they are not altered by emergency. \ldots\ While emergency does not create power, emergency may furnish the occasion for the exercise of power.''
\end{quote}
\noindent --- Chief Justice Hughes, 290 U.S. at 425--26.

\begin{quote}
``The reserved power cannot be construed so as to destroy the limitation, nor is the limitation to be construed to destroy the reserved power in its essential aspects.''
\end{quote}
\noindent --- Chief Justice Hughes, 290 U.S. at 439.

\end{document}
